\SourceLeaf{306}
\Needspace{5\baselineskip}\begin{center}
[2.]
\end{center}

Wenn $B, C$ zwei Punkte auf $a$ sind, kann man ohne Weiteres $B - C = ai$ setzen.

\Needspace{5\baselineskip}\begin{center}
[3.]
\end{center}

Die Hauptfrage ist nun, welche Kräfte man im baryzentrischen Kalkül an den Eckpunkten oder im Resultantenkalkül an den Seiten des Pentagons ansetzen muss, damit zwischen ihnen und ihren Äquivalenten am inneren Pentagon möglichst einfache Verhältnisse bestehen.

\Needspace{5\baselineskip}\begin{center}
[4.]
\end{center}

Die Kollinearität zweier Systeme aus Punkten oder Geraden lässt sich wohl am treffendsten so definieren, dass sich immer zwischen je 5 Punkten (Geraden) des einen Systems und den 5 entsprechenden des anderen dieselbe lineare Gleichung aufstellen lässt.

\Needspace{5\baselineskip}\begin{center}
[5.]
\end{center}

Das Charakteristische der Kollineation lässt sich wohl am einfachsten so ausdrücken.

Es seien $A, B.\,.\,.$ Geraden und $P,Q.\,.\,.$ Punkte, alle im selben System, denen im anderen die Geraden $a,b.\,.\,.$ und die Punkte $p,q.\,.\,.$ entsprechen. Bezeichnet man nun mit $PA$ den (senkrechten) Abstand des Punktes $P$ von der Geraden $A$, so gilt allgemein
\[
\frac{PA.QB.pb.qa}{pa.qb.PB.QA} = 1.
\]
